\documentclass[12pt]{article}
\usepackage[T1]{fontenc}
\usepackage[utf8]{inputenc}
\usepackage[brazil]{babel}
\usepackage[margin=1in]{geometry}
\usepackage{ragged2e}
\usepackage[normalem]{ulem}
\setlength{\parindent}{0pt}
\setlength{\parskip}{0.8\baselineskip}
\begin{document}
\justifying
{\Large\bfseries \textbf{PU-AC }\textbf{DISSONANCIA}\textbf{ COM TEMA TNU-}\textbf{ADEQUAÇAO}}\\

\noindent Verifica-se que a questão objeto da controvérsia jurídica suscitada no recurso manejado pela parte recorrente esteve afetada no representativo da controvérsia - \textbf{Tema n. }\textbf{9} da TNU - (PU no processo n. PEDILEF 0027714-87.2007.4.01.3600/ MT),  afetado em 06/09/2011 foi julgado (trânsito em julgado em 04/03/2013), por meio do qual foi elucidada a seguinte questão:

\noindent \textit{"Saber se os Decretos n. 5.554/2005, 5.992/2006 e 6.258/2007 reajustaram o valor das diárias dos servidores da FUNASA."}

\noindent Restou firmada a seguinte tese:

\noindent \textit{"Os Decretos n. 5.554/2005, 5.992/2006 e 6.258/2007, não reajustaram o valor das diárias dos servidores da FUNASA, apenas ampliaram o rol dos destinos a serem percorridos. Vide Tema 92 - questão similar." (Tema 9 TNU).}

\noindent Com efeito, depreende-se da leitura do voto condutor do acórdão que o mesmo se encontra em desconformidade com o decidido no  tema pela TNU.

\noindent A proposito, consta do acordao hostilizado: (...) NAO\_INFORMADO

\noindent Posto isso, tendo em vista que o Acórdão recorrido está em dissonância com o entendimento dominante do Órgão de Uniformização, com arrimo no art. 14, IV, \textit{b,} da Resolução CJF n. 586/2019, determino a \textbf{REMESSA dos autos} \textbf{à Turma de origem, }a fim de que o Colegiado proceda, se for o caso, à retratação do julgado, adequando-o aos termos acolhidos pela Instância Superior.

\noindent Após o julgamento pela Turma de Origem, observe-se o disposto nos §§7º e 8º do art. 14 da Resolução CJF n. 586/2019:

\noindent \textit{§ 7º Nos casos do inciso IV, a nova decisão proferida pela Turma de origem substitui a anterior, ficando integralmente prejudicados os pedidos de uniformização de interpretação de lei federal anteriormente interpostos.}

\noindent \textit{§ 8º Interposto novo pedido de uniformização de interpretação de lei federal em face da decisão prevista no §7º, não cabe nova remessa à Turma de origem nos termos do inciso IV, devendo se prosseguir no exame de admissibilidade.}

\noindent Intimem-se. Cumpra-se.
\end{document}
